\section*{第 \ref{ch:g9:atoms-to-galaxies} 章 --- 从原子到星系：宇宙的尺度}

\begin{solution}{exo:g9:atoms-to-galaxies:1}
原子核（$10^{-15}$ \unit{m}）、原子（$10^{-10}$）、分子
（约 $10^{-9}$）、蚂蚁（$10^{-2}$，
若取较小的蚂蚁则为 $10^{-3}$）。
\end{solution}

\begin{solution}{exo:g9:atoms-to-galaxies:2}
极小而沉重、带正电的原子核位于中心；
几乎没有重量、带负电的电子云则远在其外。
大小方面：原子约 $10^{-10}$ \unit{m}，
原子核约 $10^{-15}$ ——相差十万倍的差距，
正是造就“体育场中一粒豌豆”这一比喻的原因。
\end{solution}

\begin{solution}{exo:g9:atoms-to-galaxies:3}
电子带负电，由 $-$ 极向 $+$ 极运动，
与约定方向相反。定律依然成立，
因为它们从不取决于谁在运动：电流、节点电流之和、欧姆
定律的图线与电费的计算，
读数完全一样，箭头依全球约定而定。
\end{solution}

\begin{solution}{exo:g9:atoms-to-galaxies:4}
取最接近的 10 的幂。身高：$10^{0}$ \unit{m}；
教室：$10^{1}$；地球直径：$10^{7}$。
\end{solution}

\begin{solution}{exo:g9:atoms-to-galaxies:5}
从原子到原子核：五个梯级。从你到地球：七个梯级。
从太阳到最近恒星的距离（$10^{11}$）与（$4 \times 10^{16}$）之比：
相差五至六个梯级——约几十万倍。
\end{solution}

\begin{solution}{exo:g9:atoms-to-galaxies:6}
原子大部分是空的，直径约 $10^{-10}$ \unit{m}，
中心是 $10^{-15}$ 的原子核。把原子堆成桌子，空隙也随之而来。
固体之所以坚实，是电子云的静电斥力
阻止它们相互重叠——空虚，却被牢牢守住。
\end{solution}

\begin{solution}{exo:g9:atoms-to-galaxies:7}
毫米，$10^{-3}$；原子，$10^{-10}$：相差七个梯级——
排成一行约有一千万个原子。
\end{solution}

\begin{solution}{exo:g9:atoms-to-galaxies:8}
光年：$10^{16}$ \unit{m}（约九点五千万亿
）。银河：$10^{21}$。已观测的最深天空：
接近 $10^{26}$。
\end{solution}

\begin{solution}{exo:g9:atoms-to-galaxies:9}
$10^{5}$ 粒豌豆：$0.01 \times 10^{5} = \qty{1000}{m}$ ——
宽一千米的“体育场”：比任何真实的都宏大，但
这幅图景依然成立——中心是一粒豌豆，
相邻的豌豆相距一千米。
\end{solution}

\begin{solution}{exo:g9:atoms-to-galaxies:10}
海洋中的水滴：约 $10^{21} \times 10^{5} = 10^{26}$
——因此，若照字面理解，旧说法便夸大了：海洋中的水滴
数比一滴水中的分子数多出好几个梯级。
较诚实的说法——一滴水中的分子数，多于一生
所能数过的任何东西——仍然成立；梯级让即便
心爱的说法也保持诚实。
\end{solution}

\begin{solution}{exo:g9:atoms-to-galaxies:11}
从原子往下约十级（自 $10^{0}$ 到 $10^{-10}$），
向上直到深空约二十六级，大致居中，
略偏向微小的一端。九年的学习足以让读者
诚实地触及所测量事物的
两端。
\end{solution}

\begin{solution}{exo:g9:atoms-to-galaxies:12}
答案因人而异——这道题本身就是那封信。
（仅为示范，举例如下：致读完“影子”一章的
一年级读者——“路面上那个黑色的孪生者，
终有一天会解释日食”；致读完“能量”一章的
五年级读者——“你追踪的那些能量链，终将
变成含焦耳的公式”；致读完欧姆定律的
八年级读者——“你画的图线，正是检验
定律的习惯，
这便是全部秘诀”。）
\end{solution}

\begin{solution}{pb:g9:atoms-to-galaxies:1}
\textbf{1.} 校园：第 $2$ 秒；地球：第 $7$ 秒；
月球轨道：第 $9$ 秒。
\textbf{2.} 第 $11$ 秒（距离 $1.5 \times 10^{11}$ \unit{m}）；
旁白：“毯子上的阳光，已是八
分钟前的旧光了。”
\textbf{3.} 大约第 $13$ 至 $16$ 秒：画面中
整个太阳系只是一个光点，尚无恒星——
说实话，接近漆黑的空旷，持续数秒：这是影片
最真实的镜头。
\textbf{4.} 已观测的最深天空，充满 $10^{26}$ \unit{m} 范围内的
星系——其最古老的光比地球本身还要古老。
\textbf{5.} 蚂蚁：向内第 $2$ 秒；细胞：第 $5$ 秒；
分子：第 $9$ 秒；原子：第 $10$ 秒。
\textbf{6.} 穿过原子自身的空旷——
电子云与中心之间是五秒的虚无：“若这个原子
是一座体育场，我们已离开最高的看台，
仍在朝一粒豌豆飞去。”
\textbf{7.} 原子核——坦率地说：“我们的梯子
到此为止，并非自然止步，而是因为这些
已是人类测量过的最小梯级；大厦
仍在延伸。”
\textbf{8.} 向外 $26$ 秒，向内 $15$ 秒：
一个四十一秒的宇宙。
\textbf{9.} 每一秒都使视野放大十倍——
相等的是\emph{比值}，而非相等的加数：梯子的
逻辑是：每一级都是一个因子，按因子计，
街道级的一步与星团级的一步一样大。
\textbf{10.} 例如：“无论看上去多满，几乎
这里的一切都是空的——物质与天空同样是
广阔虚无中稀疏的孤岛。”
\textbf{11.} 显微镜（透镜一章：两块凸透镜
用于观察微小之物）与望远镜（同一章，
收集遥远的微光）——两个方向的钥匙，
都由本书的光学原理磨制而成。
\textbf{12.} 例如：“九年，一把梯子：从
瓢虫可以数清的腿，到比脚下大地还古老的星系；
每一级都由任何人都可以核对的测量确定。
下一年，我们留下的那些问题——力与
运动、光的本质、原子的核心——开始
得到解答；带着这把梯子来吧。”
\end{solution}
